\begin{lemma}[A rational upper-start Newton deficit]
\label{lem:f-rational-newton-deficit}
Let $0<m<1/2$, $h=\sqrt3/2$, and let $C_L(m)$ be the unique root in
$[h,1]$ of $F_m(c)=c^4-c^2+mc-m^2$.
Set $\omega=m(1-m)$ and $c_0=1-\omega/2$. Then
\[
 r(m):=1-c_0+\frac{F_m(c_0)}{F_m'(c_0)}
 =\frac{\omega(4-\omega)(3\omega^2-4\omega+4)}
 {8(4+2m-10\omega+6\omega^2-\omega^3)}
\]
satisfies
\[
0<\omega/2<r(m)\le1-C_L(m)<m.
\]
\end{lemma}
\begin{proof}
One has $c_0\ge7/8>h$ and
\[
 F_m(c_0)=\frac{m^2(1-m)}{16}P(m),\qquad
 P(m)=12-28m+17m^2-11m^3+3m^4-m^5.
\]
On $[0,1/2]$, $P'\le-28+17+3/2=-19/2$ and
$P(1/2)=33/32>0$. Thus $F_m(c_0)>0$.
The polynomial $F_m$ is increasing and strictly convex on $[h,1]$,
so $c_0>C_L(m)$ and its tangent zero
$c_1=c_0-F_m(c_0)/F_m'(c_0)$ satisfies
$C_L(m)\le c_1<c_0$. This proves the first bounds.
If $1-m<h$, then $C_L(m)>1-m$ immediately. Otherwise it follows from
\[
 F_m(1-m)=-m((1-m)^3+m^2)<0.
\]
The displayed rational formula follows by substitution; its denominator
factor $4+2m-10\omega+6\omega^2-\omega^3=2F_m'(c_0)$ is positive.
\end{proof}

\begin{lemma}[Uniform affine radius envelope]
\label{lem:f-uniform-affine-radius-envelope}
\label{lem:paper-branchwise-cbar}
In the strict F domain, put $U=a+b-1$, $m=\min(1-a,1-b)$, and let
$c_*$ be the exact common-pair capacity. For any
$0<r\le1-C_L(m)$, define
\[
 E_r(m,U)=r+(m-r)\left(1-\frac Ur\right)_+
 =\max\left\{r,m-\frac{m-r}{r}U\right\}.
\]
Then
\[
 1-c_*\ge E_r(m,U).
\]
In particular, the rational function in
Lemma~\ref{lem:f-rational-newton-deficit} gives one explicit radius recipe
without testing the original capacity-branch selector.
\end{lemma}
\begin{proof}
Reflect to $a\ge b$ and set $s=1-U$, $C=C_L(m)$. The exact capacity
formula is $c_*=C$ on the quartic branch and
\[
 g_m(s)=\frac{2(s-m)}{1+\sqrt{4s^2-3}}
\]
on the triangular branch. On the latter, its selector
$\chi=F_m(s)>0$ gives $s>h$ because
$\chi\le s^4-3s^2/4$, and hence $s>C$.
Writing $E=\sqrt{4s^2-3}$,
\[
 g_m'(s)=\frac{2(E-3+4ms)}{E(1+E)^2}\le0
 \qquad(C<s\le1).
\]
The root equation gives
$m/C=(1\pm\sqrt{4C^2-3})/2$, and substitution shows $g_m(C)\le C$.
Thus $c_*\le C$ on both branches, so $1-c_*\ge r$.

It remains to consider $0<U<r$. Here $s>1-r\ge C$, so the triangular
formula applies. On the analytic interval $s\in[1-r,1]$,
\[
 g_m''(s)=\frac{8\{s(3+9E-2E^2)-m(3+9E+2E^3)\}}
 {E^3(1+E)^3}>0.
\]
Indeed $m/s\le1/\sqrt3<2/3$ and
\[
 (3+9E-2E^2)-\frac23(3+9E+2E^3)
 =1+3E-2E^2-\frac43E^3\ge\frac23
 \quad(0\le E\le1),
\]
since the last cubic is concave and lies above its endpoint chord.
Endpoint cases follow by continuity. Therefore $U\mapsto g_m(1-U)$
is convex. Its endpoint bounds are
$g_m(1)=1-m$ and $g_m(1-r)\le C\le1-r$.
The chord inequality gives
\[
 c_*\le(1-U/r)(1-m)+(U/r)(1-r),
\]
which is the claimed affine bound. Combining it with $1-c_*\ge r$
proves the result.
\end{proof}
